\documentclass[12pt]{report}
\usepackage[brazilian]{babel}
\usepackage{csquotes}
\usepackage{makeidx}
\usepackage{glossaries}
\usepackage[
  perfil=academico,
  tipo-trabalho=dissertacao,
  impressao=frente-e-verso,
  citacoes=autor-data
]{abntex3}

\makeindex
\makeglossaries
\newglossaryentry{repositorio}{
  name={repositório digital},
  description={sistema para armazenar e disponibilizar objetos digitais}
}
\addbibresource{referencias-exemplo.bib}
\ABNTEXsetup{
  titulo={Modelo canônico completo de dissertação},
  subtitulo={base auditável com todos os elementos implementados},
  autoria={Maria da Silva},
  instituicao={Universidade Exemplo},
  natureza={Dissertação},
  objetivo={apresentada para obtenção do título de Mestra},
  area={Ciência da Informação},
  orientacao={Prof. Dr. João Souza},
  coorientacao={Profa. Dra. Ana Santos},
  local={Recife},
  data={2026}
}
\ABNTEXabreviatura{org.}{organização}
\ABNTEXsigla{ABNT}{Associação Brasileira de Normas Técnicas}
\ABNTEXsimbolo{$n$}{Tamanho da amostra}

\begin{document}
\ABNTEXacademicoexterna
\ABNTEXlombada[Dissertação 2026][Universidade Exemplo]

\ABNTEXacademicopretextual
\ABNTEXfolhaderosto
\ABNTEXfichacatalografica{%
  Ponto de extensão para dados catalográficos preparados por profissional.
}
\ABNTEXerrata
  {SILVA, Maria da. \textbf{Modelo canônico completo de dissertação}.
   2026. Dissertação — Universidade Exemplo, Recife, 2026.}
  {\begin{center}
     \begin{tabular}{llll}
       Folha & Linha & Onde se lê & Leia-se\\
       12 & 3 & repositorio & repositório
     \end{tabular}
   \end{center}}
\ABNTEXfolhadeaprovacao
  {30 de julho de 2026}
  {%
    \ABNTEXmembrobanca
      {Prof. Dr. João Souza}{Orientador}{Universidade Exemplo}
    \ABNTEXmembrobanca
      {Profa. Dra. Ana Santos}{Coorientadora}{Universidade Exemplo}
    \ABNTEXmembrobanca
      {Profa. Dra. Beatriz Alves}{Examinadora}{Instituto Exemplo}
  }
\ABNTEXdedicatoria{À comunidade que mantém o conhecimento acessível.}
\ABNTEXagradecimentos{À equipe, à instituição e às pessoas avaliadoras.}
\ABNTEXepigrafe[Autoria demonstrativa]{Documentar é tornar verificável.}
\ABNTEXresumo
  {Modelo completo de dissertação para reutilização e auditoria dos elementos
   atualmente implementados pelo pacote.}
  {dissertação, documentação, auditoria, LaTeX}
\ABNTEXresumoemlinguaestrangeira
  {Complete dissertation template for reuse and auditing of every element
   currently implemented by the package.}
  {dissertation, documentation, audit, LaTeX}
\ABNTEXlistadeilustracoes
\ABNTEXlistadetabelas
\ABNTEXlistadeabreviaturasesiglas
\ABNTEXlistadesimbolos
\ABNTEXsumario

\ABNTEXacademicotextual
\chapter{Introdução}
\index{modelo}
Apresentam-se tema, problema, hipótese, objetivos e justificativa. O
\gls{repositorio} é o objeto demonstrativo. A fundamentação bibliográfica é
processada externamente \parencite{associacao2025}.
\section{Problema e hipótese}
O problema investiga a auditabilidade; a hipótese relaciona documentação e
reprodutibilidade.
\section{Objetivos e justificativa}
O objetivo é exercitar a API completa e a justificativa é oferecer uma base
segura para novos documentos.

\chapter{Referencial teórico}
\section{Apresentação documental}
\textcite{associacao2025} exemplifica uma chamada integrada.
\begin{ABNTEXcitacaolonga}
Trecho demonstrativo para inspeção do recuo, corpo e espaçamento destinados
a citações longas, sem reproduzir conteúdo de norma técnica.
\end{ABNTEXcitacaolonga}
\ABNTEXnota{Nota de rodapé demonstrativa.}

\chapter{Metodologia}
\section{Coleta e análise}
O procedimento inventaria comandos, compila o documento e registra resultados.

\chapter{Resultados e discussão}
\begin{figure}[ht]
  \centering
  \rule{0.55\linewidth}{2cm}
  \caption{Representação do fluxo de validação}
  \ABNTEXfonte{elaborada pela autora.}
\end{figure}
\begin{table}[ht]
  \centering
  \caption{Resultados demonstrativos}
  \begin{tabular}{ll}
    Verificação & Resultado\\
    Compilação & concluída\\
    Bibliografia & processada
  \end{tabular}
  \ABNTEXfonte{elaborada pela autora.}
\end{table}

\chapter{Conclusão}
Os objetivos são retomados e as limitações do modelo são explicitadas.

\ABNTEXacademicopostextual
\ABNTEXbibliografia
\ABNTEXglossario
\ABNTEXappendix{Formulário elaborado pela autora}
Instrumento produzido durante a pesquisa.
\ABNTEXannex{Documento institucional de demonstração}
Material externo associado à pesquisa.
\ABNTEXindice[Índice de assuntos]
\end{document}
